\section*{Hoofdstuk \ref{ch:g11:polarity-and-cohesion} --- Elektronegativiteit, polariteit en krachten tussen moleculen}

\begin{solution}{exo:g11:polarity-and-cohesion:1}
\ce{H-Cl}: chloor (3.16 tegen 2.20). \ce{O-H}: zuurstof (3.44).
\ce{C-O}: zuurstof (3.44 tegen 2.55).
\end{solution}

\begin{solution}{exo:g11:polarity-and-cohesion:2}
\ce{H-H}: 0, niet polair. \ce{C-H}: 0.35, geldt als apolair. \ce{O-H}:
1.24, polair. \ce{N-H}: 0.84, polair. \ce{Cl-Cl}: 0, niet polair.
\ce{C-Cl}: 0.61, polair.
\end{solution}

\begin{solution}{exo:g11:polarity-and-cohesion:3}
Stikstof (3.04 > 2.19) en zuurstof (3.44 > 2.58): de \omterm{def:g11:polarity-and-cohesion:electronegativity}{elektronegativiteit}
neemt af naar beneden in een groep. Stikstof (3.04 > 2.55) en magnesium
(1.31 > 0.93): ze neemt toe van links naar rechts langs een periode.
\end{solution}

\begin{solution}{exo:g11:polarity-and-cohesion:4}
Ammoniak en water: elk heeft waterstofatomen gebonden aan stikstof of
zuurstof, en \omterm{def:g10:lewis-and-shape:lone-pair}{vrije elektronenparen} op dat \omterm{def:g7:atoms-and-molecules:atom}{atoom}. Methaan heeft geen
waterstof aan N, O of F; waterstofchloride heeft waterstof gebonden aan
chloor, en dat is niet een van de drie \omterm{def:g7:atoms-and-molecules:atom}{atomen} uit de definitie.
\end{solution}

\begin{solution}{exo:g11:polarity-and-cohesion:5}
Ja: \omterm{def:g11:polarity-and-cohesion:van-der-waals}{vanderwaalsbindingen} bestaan tussen alle \omterm{def:g7:atoms-and-molecules:molecule}{moleculen}. Maar
methaanmoleculen zijn klein (10 \omterm{def:g9:inside-the-atom:nucleus}{elektronen}), dus deze bindingen zijn heel
zwak en breken al bij een heel lage temperatuur: methaan kookt bij
ongeveer \qty{-162}{\celsius}.
\end{solution}

\begin{solution}{exo:g11:polarity-and-cohesion:6}
\ce{CH2Cl2}: polair (twee polaire \ce{C-Cl}-bindingen en twee bijna
apolaire \ce{C-H}-bindingen: geen symmetrie die ze opheft). \ce{CCl4}:
apolair (vier gelijke bindingen op de hoekpunten van een tetraëder).
\ce{NH3}: polair (piramide). \ce{CO2}: apolair (lineair, tegengestelde
bindingen). \ce{BF3}: apolair (drie gelijke \omterm{def:g11:polarity-and-cohesion:polar-bond}{polaire bindingen} onder
\ang{120} in een vlak heffen elkaar op).
\end{solution}

\begin{solution}{exo:g11:polarity-and-cohesion:7}
De waterstof van de \ce{O-H}-groep kan gegeven worden (niet die van
\ce{CH3}, die aan koolstof gebonden zijn); het zuurstofatoom, met zijn
twee \omterm{def:g10:lewis-and-shape:lone-pair}{vrije elektronenparen}, kan opnemen. Ja: de \ce{O-H} van methanol kan
binden aan de zuurstof van een watermolecuul, en een \ce{O-H} van water
aan de zuurstof van methanol, bijvoorbeeld
\ce{CH3-O-H}\,$\cdots$\,\ce{OH2}.
\end{solution}

\begin{solution}{exo:g11:polarity-and-cohesion:8}
De drie \omterm{def:g7:atoms-and-molecules:molecule}{moleculen} zijn apolair en hebben dezelfde vorm, maar zijn steeds
groter (10, 18 en 36 \omterm{def:g9:inside-the-atom:nucleus}{elektronen}): hun \omterm{def:g11:polarity-and-cohesion:van-der-waals}{vanderwaalsbindingen} worden
sterker, en hun kookpunten stijgen mee.
\end{solution}

\begin{solution}{exo:g11:polarity-and-cohesion:9}
\omterm{def:g9:periodic-table-first-look:period}{Groep} 14: methaan heeft geen waterstof gebonden aan N, O of F en vormt
geen \omterm{def:g11:polarity-and-cohesion:hydrogen-bond}{waterstofbruggen}, dus het volgt de trend van zijn groep.
\end{solution}

\begin{solution}{exo:g11:polarity-and-cohesion:10}
De \omterm{def:g11:polarity-and-cohesion:van-der-waals}{vanderwaalsbindingen}: de \omterm{def:g7:atoms-and-molecules:molecule}{moleculen} worden groter van \ce{HCl} naar
\ce{HI}, en de toename van deze bindingen met de grootte weegt zwaarder
dan de afname van de aantrekking tussen de \omterm{def:g11:polarity-and-cohesion:polar-bond}{partiële ladingen}.
\end{solution}

\begin{solution}{exo:g11:polarity-and-cohesion:11}
Een joodmolecuul is veel groter dan een chloormolecuul (106 \omterm{def:g9:inside-the-atom:nucleus}{elektronen}
tegen 34): zijn \omterm{def:g11:polarity-and-cohesion:van-der-waals}{vanderwaalsbindingen} zijn veel sterker, sterk genoeg om
de \omterm{def:g7:atoms-and-molecules:molecule}{moleculen} bij kamertemperatuur in een vaste stof bij elkaar te
houden.
\end{solution}

\begin{solution}{exo:g11:polarity-and-cohesion:12}
Propaan is apolair: alleen \omterm{def:g11:polarity-and-cohesion:van-der-waals}{vanderwaalsbindingen}. Methoxymethaan is polair
(twee polaire \ce{C-O}-bindingen in een gebogen \ce{C-O-C}) maar heeft
geen waterstof aan zijn zuurstof: geen \omterm{def:g11:polarity-and-cohesion:hydrogen-bond}{waterstofbruggen} tussen zijn
\omterm{def:g7:atoms-and-molecules:molecule}{moleculen}. Ethanol is polair en zijn \ce{O-H}-groepen vormen
\omterm{def:g11:polarity-and-cohesion:hydrogen-bond}{waterstofbruggen}. Hoe sterker de aantrekking, hoe hoger het kookpunt:
propaan < methoxymethaan < ethanol.
\end{solution}

\begin{solution}{exo:g11:polarity-and-cohesion:13}
Het open netwerk van ijs neemt meer ruimte in dan dezelfde \omterm{def:g7:atoms-and-molecules:molecule}{moleculen} in
de vloeistof, dus ijs is minder dicht dan vloeibaar water en drijft. Een
meer bevriest van bovenaf: de ijslaag isoleert het water eronder, dat
vloeibaar blijft, en de vissen overleven.
\end{solution}

\begin{solution}{exo:g11:polarity-and-cohesion:14}
Vast methaan (\omterm{def:g11:polarity-and-cohesion:van-der-waals}{vanderwaalsbindingen} tussen kleine \omterm{def:g7:atoms-and-molecules:molecule}{moleculen}) < ijs
(\omterm{def:g11:polarity-and-cohesion:hydrogen-bond}{waterstofbruggen}) < kaliumchloride (aantrekking tussen \omterm{def:g9:ions:ion}{ionen}).
\end{solution}

\begin{solution}{exo:g11:polarity-and-cohesion:15}
Een watermolecuul heeft twee waterstofatomen om te geven en twee \omterm{def:g10:lewis-and-shape:lone-pair}{vrije
elektronenparen} om mee op te nemen: elk \omterm{def:g7:atoms-and-molecules:molecule}{molecuul} kan aan vier
\omterm{def:g11:polarity-and-cohesion:hydrogen-bond}{waterstofbruggen} meedoen, twee gegeven en twee opgenomen, wat een
volledig netwerk opbouwt. Een \omterm{def:g7:atoms-and-molecules:molecule}{molecuul} waterstoffluoride heeft drie \omterm{def:g10:lewis-and-shape:lone-pair}{vrije
elektronenparen} maar maar één waterstofatoom: het geeft maar één
\omterm{def:g11:polarity-and-cohesion:hydrogen-bond}{waterstofbrug}, dus kan er gemiddeld maar één brug per \omterm{def:g7:atoms-and-molecules:molecule}{molecuul} ontstaan
(zigzagketens). Water heeft ongeveer twee keer zoveel \omterm{def:g11:polarity-and-cohesion:hydrogen-bond}{waterstofbruggen}
per \omterm{def:g7:atoms-and-molecules:molecule}{molecuul}, en kookt hoger.
\end{solution}

\begin{solution}{pb:g11:polarity-and-cohesion:1}
\textbf{1.} \ce{H-O-H}, \ce{H-S-H}, \ce{H-Se-H}, elk centraal \omterm{def:g7:atoms-and-molecules:atom}{atoom} met
twee \omterm{def:g10:lewis-and-shape:lone-pair}{vrije elektronenparen}: O, S en Se hebben zes \omterm{def:g10:electron-shells:valence}{valentie-elektronen},
omdat ze in dezelfde groep staan.

\textbf{2.} $M(\ce{H2O}) = \qty{18.0}{g/mol}$, $M(\ce{H2S}) =
\qty{34.1}{g/mol}$, $M(\ce{H2Se}) = \qty{81.0}{g/mol}$.

\textbf{3.} $3.44 - 2.20 = 1.24$; $2.58 - 2.20 = 0.38$;
$2.55 - 2.20 = 0.35$. Alleen de \ce{O-H}-binding is duidelijk polair.

\textbf{4.} Gebogen (vier groepen, waarvan twee vrije paren). Water is
polair; waterstofsulfide en waterstofselenide maar heel weinig.

\textbf{5.} 10, 18 en 36 \omterm{def:g9:inside-the-atom:nucleus}{elektronen}: waterstofselenide heeft de sterkste
\omterm{def:g11:polarity-and-cohesion:van-der-waals}{vanderwaalsbindingen}.

\textbf{6.} Alleen water: zijn waterstofatomen zijn gebonden aan
zuurstof, een van de drie heel elektronegatieve \omterm{def:g7:atoms-and-molecules:atom}{atomen}; zwavel en seleen
zijn niet elektronegatief genoeg om hun waterstofatomen een duidelijke
$\delta^+$ te geven.

\textbf{7.} Vier: twee via zijn eigen waterstofatomen, twee opgenomen
door zijn twee \omterm{def:g10:lewis-and-shape:lone-pair}{vrije elektronenparen}.

\textbf{8.} \omterm{def:g11:polarity-and-cohesion:van-der-waals}{Vanderwaalsbindingen} (met daarbij een zwakke aantrekking
tussen de kleine \omterm{def:g11:polarity-and-cohesion:polar-bond}{partiële ladingen}).

\textbf{9.} $212.9 - 273.15 = \qty{-60.3}{\celsius}$ en
\qty{-41.3}{\celsius}.

\textbf{10.} $-41.3 - (-60.3) = \qty{19.0}{\celsius}$.

\textbf{11.} Het \omterm{def:g7:atoms-and-molecules:molecule}{molecuul} is groter, dus zijn \omterm{def:g11:polarity-and-cohesion:van-der-waals}{vanderwaalsbindingen} zijn
sterker.

\textbf{12.} $-60.3 - 19.0 = \qty{-79.3}{\celsius}$.

\textbf{13.} Stijging $-88.1 - (-112) = \qty{23.9}{\celsius}$; voorspelling
voor methaan $-112 - 23.9 = \qty{-136}{\celsius}$, tegen
\qty{-162}{\celsius} in werkelijkheid: de extrapolatie ligt ongeveer
\qty{26}{\celsius} te hoog. Ze is niet exact; de echte trend is geen
rechte lijn.

\textbf{14.} Stijging $\qty{18.7}{\celsius}$; voorspelling voor
waterstoffluoride $-85.1 - 18.7 = \qty{-103.8}{\celsius}$, tegen
\qty{19.6}{\celsius}: een verschil van ongeveer \qty{123}{\celsius}.

\textbf{15.} $100 - (-79.3) = \qty{179}{\celsius}$, ongeveer
\qty{180}{\celsius}.

\textbf{16.} Stijging $\qty{25.3}{\celsius}$; voorspelling $-87.8 - 25.3 =
\qty{-113}{\celsius}$; verschil $-33.4 - (-113) = \qty{80}{\celsius}$.

\textbf{17.} Water (\qty{180}{\celsius}) > waterstoffluoride
(\qty{123}{\celsius}) > ammoniak (\qty{80}{\celsius}). Water geeft twee
waterstofatomen en neemt op met twee \omterm{def:g10:lewis-and-shape:lone-pair}{vrije elektronenparen}: vier
\omterm{def:g11:polarity-and-cohesion:hydrogen-bond}{waterstofbruggen} per \omterm{def:g7:atoms-and-molecules:molecule}{molecuul}, allemaal benut. Waterstoffluoride heeft
één waterstofatoom om te geven, en ammoniak maar één \omterm{def:g10:lewis-and-shape:lone-pair}{vrij elektronenpaar}
om mee op te nemen: gemiddeld minder bruggen per \omterm{def:g7:atoms-and-molecules:molecule}{molecuul}, en stikstof,
dat minder elektronegatief is, vormt zwakkere.

\textbf{18.} Een gas: het zou koken bij ongeveer \qty{-80}{\celsius}. De
aarde zou geen vloeibaar water hebben, geen oceanen, geen regen, en geen
van het leven dat daarvan afhangt.

\textbf{19.} Het extrapoleert een rechte lijn door maar twee punten, en de
test op \omterm{def:g9:periodic-table-first-look:period}{groep} 14 laat zien dat zo'n extrapolatie zo'n
\qty{25}{\celsius} ernaast kan zitten.

\textbf{20.} Zonder \omterm{def:g11:polarity-and-cohesion:hydrogen-bond}{waterstofbruggen} zou water koken bij ongeveer
\qty{-80}{\celsius}, zo'n \qty{180}{\celsius} onder zijn echte kookpunt.
\end{solution}
